% from ega4-13-fr.tex
\oldpage[IV-3]{199}
considérée comme sous-préschéma fermé réduit de $Y$). Par exemple, soient $k$
un corps, $A_0=k[S,T]$ l'anneau de polynômes à deux indéterminées,
$\mathfrak a=\mathfrak p_1\mathfrak p_2$, où $\mathfrak p_1$ et
$\mathfrak p_2$ sont les idéaux premiers $A_0S$ et $A_0T$ de $A_0$~; soient
$A=A_0/\mathfrak a$ et $Y=\operatorname{Spec}(A)$, qui a deux composantes
irréductibles $Y_1=\operatorname{Spec}(A/\mathfrak p_1)$,
$Y_2=\operatorname{Spec}(A/\mathfrak p_2)$~; prenons $X=Y_1$,
$f:X\to Y$ étant l'injection canonique, qui est évidemment un morphisme
équidimensionnel. Prenons d'autre part $Y'=Y_2$, $g:Y'\to Y$ étant l'injection
canonique~; alors on a
$X'=X\times_Y Y'=\operatorname{Spec}((A/\mathfrak p_1)\otimes_A
(A/\mathfrak p_2))=\operatorname{Spec}(A/(\mathfrak p_1+\mathfrak p_2))$,
et $\mathfrak p_1+\mathfrak p_2$ est un idéal maximal de $A$, donc $X'$ est
réduit à un point, et le morphisme $f':X'\to Y'$ n'est \emph{pas dominant},
l'image par $f'$ de l'unique point de $X'$ étant un point fermé de $Y'$~;
donc $f'$ n'est pas équidimensionnel.

On peut aussi donner un contre-exemple où $X$ et $Y$ sont \emph{intègres}, $f$
\emph{fini} et \emph{birationnel} (et \emph{a fortiori} équidimensionnel par
(13.3.1, $b$)), $g:Y'\to Y$ \emph{fini et dominant}. Soient $A$ et
$\overline A$ les anneaux locaux définis dans (11.7.5), et prenons
$Y=\operatorname{Spec}(A)$, $X=\operatorname{Spec}(\overline A)$~; d'autre
part, avec les notations de (11.7.5), prenons
$Y'=\operatorname{Spec}(\overline B)$~; alors
$X'=\operatorname{Spec}(\overline A\otimes_A\overline B)
=\operatorname{Spec}(\overline B\otimes_B\overline B)$~; mais on vérifie
aussitôt que $\overline B\otimes_B\overline B$ est composé direct des anneaux
$B/\mathfrak p'$, $B/\mathfrak p''$ et de deux anneaux isomorphes à
$B/\mathfrak n$, dont les spectres sont donc réduits à un point~; comme les
projections de ces points sont des points fermés de $Y'$, ici encore on voit
que $f'$ ne transforme pas une composante irréductible de $X'$ en une partie
partout dense d'une composante irréductible de $Y'$, donc $f'$ n'est pas
équidimensionnel.

Dans cet exemple, l'anneau $A$ n'est pas géométriquement unibranche~; nous
allons voir au paragraphe suivant (14.4.6) que de tels phénomènes ne peuvent se
produire lorsque les points de $Y$ sont \emph{géométriquement unibranches}. Le
manque de stabilité de la notion de morphisme équidimensionnel restreint
grandement son intérêt, au profit de la notion de morphisme universellement
ouvert, qui va être étudiée en détail au paragraphe suivant.
\end{remark}

% IV3_B28_P199_END_BEFORE_SECTION_14

% (continued in ega4-14-fr.tex, before its first page marker)
\section{Morphismes universellement ouverts}
\label{section:IV.14-fr}

Les \S\S~14 et 15 sont consacrés à l'étude de la notion de morphisme
\emph{universellement ouvert} (2.4.2). On a déjà vu (2.4.6) qu'un morphisme
\emph{plat} et localement de présentation finie est universellement ouvert, la
réciproque étant inexacte. Au \S~14, on examine d'abord les propriétés des
dimensions des fibres d'un morphisme universellement ouvert $f:X\to Y$~;
lorsque $X$ et $Y$ sont localement noethériens, $f$ se comporte à cet égard
(14.2.1) comme un morphisme plat (cf. 6.1.2), et est en particulier
\emph{équidimensionnel} lorsqu'il est localement de type fini et dominant et que
$X$ est irréductible (14.2.2). Inversement, un morphisme équidimensionnel
localement de type fini $f:X\to Y$ est universellement ouvert lorsqu'on suppose
en outre que $Y$ est \emph{géométriquement unibranche}, et en particulier lorsque
$Y$ est \emph{normal} (critère de Chevalley, 14.4.4). Nous montrons aussi que les
morphismes universellement ouverts localement de type fini $f:X\to Y$ (lorsque
$X$ et $Y$ sont
